\section{Noise is filtered, not accumulated}
\label{sec:theory}

Fix a policy $\theta_0$ and start two runs from it. They share parameters, optimiser state, data
distribution and recipe, and independently draw everything the seed controls: which prompts each
batch contains and which responses are sampled. That is the comparison a practitioner makes when
they rerun, and it is the estimand throughout, so both terms of \eqref{eq:two-level} contribute.

Write $\delta_t := \theta_t - \bar\theta_t$ for one run's departure from the mean trajectory. With
updates $\theta \leftarrow \theta + \eta\ghat$ and
$\mathcal{J}_t := \partial\gbar/\partial\theta$,
\begin{equation}
  \label{eq:seed-recursion}
  \delta_{t+1} = A_t \delta_t + \eta\,\xi_t,
  \qquad A_t := I + \eta \mathcal{J}_t,
\end{equation}
to first order in $\delta$, where $\E[\xi_t \mid \mathcal{F}_t] = 0$ and
$\Cov(\xi_t) = \Sigma_t/P$.

If the objective were flat, $\mathcal{J} = 0$, this would be a random walk: divergence would
accumulate as $\E\|\Delta_t\|^2 \propto t$ and a longer run would be strictly less reproducible.
That is the implicit model behind measuring run-to-run variance by brute force. But two RLVR runs
climb the same objective, and $A_t$ reweights a perturbation rather than passing it through.

\begin{theorem}[Finite-time covariance]
  \label{thm:propagation}
  Let $S_t := \Cov(\delta_t)$, with $S_0 = 0$ when runs start together. Under
  \eqref{eq:seed-recursion},
  \begin{equation}
    \label{eq:lyapunov}
    S_{t+1} = A_t S_t A_t^{\!\top} + Q_t,
    \qquad Q_t := \eta^2 \Sigma_t / P,
  \end{equation}
  hence $S_T = \sum_{t<T} \Phi_{T,t+1} Q_t \Phi_{T,t+1}^{\!\top}$ with
  $\Phi_{T,s} := \prod_{u=s}^{T-1} A_u$ taken in order of decreasing $u$. Two independent seeds
  satisfy $\Cov(\theta^A_T - \theta^B_T) = 2S_T$.
\end{theorem}

\begin{proof}[Proof sketch]
  Taking the covariance of \eqref{eq:seed-recursion} leaves a cross term
  $\E[\delta_t \xi_t^{\!\top}]$. It vanishes without assuming $\xi_t$ and $\delta_t$ are
  independent, which they are not, since the batch noise is drawn at $\theta_t$ and that parameter
  was produced by earlier noise. What holds is that $\xi_t$ is a martingale difference,
  $\E[\xi_t \mid \mathcal{F}_t] = 0$, and $\delta_t$ is $\mathcal{F}_t$-measurable, so the cross
  term is zero by the tower rule. Unrolling from $S_0 = 0$ gives the sum.
\end{proof}

The factor of two belongs to the difference of two runs and not to the covariance of one; keeping
them apart is what makes the statement unambiguous.

\subsection{Carrying it to the number that gets reported}

Equation \eqref{eq:lyapunov} concerns policies. What a paper reports is a score $M(\theta)$:
pass rate, exact match, win rate. Linearising about the endpoint,
$\Var_{\text{seed}}[M(\theta_T)] \approx \nabla M^{\!\top} S_T \nabla M$, which needs a matrix
nobody can form for a real model. It is not needed.

\begin{theorem}[Adjoint form of the reported spread]
  \label{thm:adjoint}
  Define $b_T := \nabla M(\theta_T)$ and $b_s := A_s^{\!\top} b_{s+1}$ for $s = T-1,\dots,0$. Then
  \begin{equation}
    \label{eq:adjoint}
      V^{\text{lin}}_T \,=\, \sum_{t=0}^{T-1} b_{t+1}^{\!\top} Q_t\, b_{t+1}
      \,=\, \frac{\eta^2}{P}\sum_{t=0}^{T-1} \Var_{i}\!\big(b_{t+1}^{\!\top} z_{t,i}\big),
  \end{equation}
  the variance taken over prompts.
\end{theorem}

\begin{proof}[Proof sketch]
  $b_s = \Phi_{T,s}^{\!\top}\nabla M$ by induction; substitute the unrolled $S_T$. The second
  equality is $Q_t = \eta^2\Sigma_t/P$ with
  $b^{\!\top}\Sigma_t b = \Var_i(b^{\!\top} z_{t,i})$, since $\Sigma_t$ is the covariance of one
  prompt's contribution rather than of the batch mean.
\end{proof}

The equality is exact for $V^{\text{lin}}_T$, the variance of the metric linearised at
$\bar\theta_T$. The quantity a practitioner wants is $\Var_{\text{seed}}[M(\theta_T)]$, and the two
agree to the order at which the linearisation holds; Section~\ref{sec:real} is about when that
agreement fails. Every term of \eqref{eq:adjoint} is one scalar per update, which a trainer can
accumulate alongside the update it is already making.

What the backward pass costs depends on $A_s^{\!\top}$, and here the estimator family pays for
itself.

\begin{proposition}[The mean update is a gradient field]
  \label{prop:potential}
  For any count-based estimator, $\gbar(\theta) = \nabla_\theta\,\Psi(\theta)$ with
  $\Psi(\theta) = \E_x[\Lambda(p_x)]$ and $\Lambda' = \lambda$, where the per-prompt mean update is
  $\lambda(p_x, G)\nabla_\theta p_x$. Hence $\mathcal{J} = \nabla^2\Psi$ is symmetric and
  $A_s^{\!\top} = A_s$.
\end{proposition}

Because $\lambda$ depends on $\theta$ only through $p_x$, $\lambda(p_x)\nabla p_x = \nabla
\Lambda(p_x)$, and averaging over $x$ commutes with the gradient. We call $\Psi$ the
\emph{estimator potential}; it is not the objective $J$ unless $\lambda \equiv 1$, and the gap is
exactly the difficulty reweighting the estimator applies. For RLOO the potential is the pass rate
itself. The practical consequence is that the backward pass needs no transpose and can be run with
the same code path as the forward one.

\subsection{Linearising the model instead of the run}
\label{sec:tangent}

Theorem~\ref{thm:adjoint} is exact in the model and first-order in the distance between runs, and
it inherits its one run's blind spot: at a run that has saturated, $\nabla M(\theta_T) = 0$ and the
forecast vanishes, wherever the other seeds went. The opposite linearisation is available. With
$\score_0(x,y) := \nabla_\theta\log\pi_{\theta_0}(y \mid x)$ the score at the starting parameters,
the \emph{tangent policy}
\begin{equation}
  \label{eq:tangent}
  \pi'_{\vartheta}(y \mid x) \;\propto\;
  \pi_{\theta_0}(y \mid x)\,\exp\!\big(\score_0(x,y)^{\!\top}\vartheta\big)
\end{equation}
is an exponential family in the frozen scores. Updating $\vartheta$ with an estimator of
\eqref{eq:family} needs sampled responses, their rewards and inner products of frozen scores, and
no pass through the model, so seeds can be \emph{simulated} by the thousand with the sampling, the
estimator's weights and the saturation all kept exactly.

\begin{proposition}[When a simulated seed is a real one]
  \label{prop:tangent}
  If $\log\policy(y \mid x)$ is affine in $\theta$ up to a per-prompt normaliser, the run from
  $\theta_0$ and the tangent run from $\vartheta_0 = 0$ have the same law, so $M(\theta_T)$ has the
  distribution of its simulated counterpart.
\end{proposition}

Both are then Markov chains with one transition kernel. A softmax over one decision qualifies. A
network does not: there a simulated run is exact in its first update and drifts as the scores
move. This is the tangent-kernel approximation \citep{jacot2018ntk, malladi2023kernel}, used to
draw many runs instead of predicting one.
